%% ========================================================================
%% Kata Pengantar
%% ========================================================================

\chapter*{Kata Pengantar}
\addcontentsline{toc}{chapter}{Kata Pengantar}

Puji syukur kehadirat Tuhan Yang Maha Esa atas segala rahmat dan karunia-Nya
sehingga buku \textit{Mastering Linux: Pemahaman Sistem Operasi untuk Pemula}
ini dapat diselesaikan.

Sistem operasi merupakan fondasi dari semua sistem komputer modern. Memahami
cara kerja sistem operasi, khususnya Linux, adalah keterampilan esensial bagi
setiap profesional teknologi informasi. Buku ini dirancang khusus untuk mendukung
pembelajaran mata kuliah Sistem Operasi dengan pendekatan praktis dan komprehensif.

\section*{Tujuan Buku}

Buku ini bertujuan untuk:

\begin{itemize}
    \item Memberikan pemahaman mendalam tentang konsep fundamental sistem operasi
    \item Mengajarkan keterampilan praktis dalam menggunakan dan mengelola sistem Linux
    \item Mempersiapkan pembaca untuk karir di bidang administrasi sistem dan DevOps
    \item Menjembatani teori akademis dengan praktik industri
\end{itemize}

\section*{Struktur Buku}

Buku ini disusun mengikuti rencana pembelajaran 16 minggu yang terstruktur:

\begin{description}
    \item[Bagian I: Fundamental Sistem Operasi]: Membahas konsep dasar, instalasi,
          dan perintah dasar Linux.
    \item[Bagian II: Manajemen File dan Proses]: Menjelaskan struktur direktori,
          operasi file, dan manajemen proses.
    \item[Bagian III: Bash Shell dan Scripting]: Mengajarkan penggunaan Bash
          dan pemrograman shell script.
    \item[Bagian IV: Manajemen Sistem Lanjutan]: Mencakup manajemen memori, user,
          layanan, aplikasi, dan backup sistem.
\end{description}

\section*{Fitur Khusus}

Buku ini dilengkapi dengan berbagai fitur pembelajaran:

\begin{itemize}
    \item \textbf{Contoh Praktis}: Setiap konsep disertai dengan contoh kode dan
          skenario nyata
    \item \textbf{Latihan}: Soal latihan di setiap akhir bab untuk memperdalam
          pemahaman
    \item \textbf{Tips dan Catatan}: Informasi tambahan dan best practices dari
          pengalaman praktisi
    \item \textbf{Referensi Perintah}: Lampiran berisi referensi lengkap perintah Linux
\end{itemize}

\section*{Ucapan Terima Kasih}

Penulis mengucapkan terima kasih kepada:

\begin{itemize}
    \item Keluarga yang telah memberikan dukungan dan kesabaran
    \item Rekan-rekan dosen dan kolega yang memberikan masukan berharga
    \item Para pembaca yang telah menjadi inspirasi dalam penulisan buku ini
    \item Tim penerbit yang telah membantu proses penerbitan
\end{itemize}

Penulis menyadari bahwa buku ini masih jauh dari sempurna. Kritik dan saran yang
membangun sangat diharapkan untuk perbaikan di edisi mendatang.

\vspace{1cm}

\noindent
Malang, April \the\year

\vspace{1cm}

\noindent
Tim Penulis,

\vspace{2cm}

\cleardoublepage
